\documentclass[11pt]{article}
\usepackage[T1]{fontenc}
\usepackage[utf8]{inputenc}
\usepackage{lmodern}
\usepackage{amsmath,amssymb,amsthm,mathtools}
\usepackage[margin=1in]{geometry}
\usepackage{microtype}
\usepackage[colorlinks=true,linkcolor=blue,citecolor=blue,urlcolor=blue]{hyperref}
\hypersetup{
 pdftitle={An interval semigroup criterion and positivity of polynomial Gaussian binomial quotients},
 pdfauthor={Zizhou Zhang},
 pdfsubject={Gaussian binomial quotients and additive semigroups},
 pdfkeywords={Gaussian binomial coefficients, coefficient positivity, additive semigroups, Kneser's theorem}
}

\newtheorem{theorem}{Theorem}[section]
\newtheorem{lemma}[theorem]{Lemma}
\newtheorem{corollary}[theorem]{Corollary}
\newtheorem{proposition}[theorem]{Proposition}
\theoremstyle{definition}
\newtheorem{definition}[theorem]{Definition}
\theoremstyle{remark}
\newtheorem{remark}[theorem]{Remark}

\newcommand{\NN}{\mathbb Z_{\geq 0}}
\newcommand{\ZZ}{\mathbb Z}
\newcommand{\CC}{\mathbb C}
\DeclareMathOperator{\diam}{diam}
\DeclareMathOperator{\stab}{stab}
\newcommand{\iv}[2]{[#1,#2]_{\ZZ}}
\newcommand{\Ga}{\Gamma}
\newcommand{\remg}[1]{\left[#1\right]_g}
\newcommand{\qbinom}[2]{\genfrac{[}{]}{0pt}{}{#1}{#2}}

\title{An interval semigroup criterion and positivity of\newline
polynomial Gaussian binomial quotients}
\author{Zizhou Zhang\\[0.5em]
\small Academy of Mathematics and Systems Science\\
\small Chinese Academy of Sciences\\
\small Beijing, China}
\date{November 26, 2026}

\begin{document}
\maketitle

\begin{abstract}
We prove that every polynomial quotient of two Gaussian binomial
coefficients with the same upper argument has nonnegative coefficients,
as conjectured by Gatzweiler and Krattenthaler. The main result is an
interval semigroup theorem. If a finite set $S$ of distinct positive
integers lies below an integer interval $A$ whose length exceeds
$\operatorname{diam}(S)$, and $A$ contains at least as many multiples of
each positive integer as $S$ does, then $A$ contains at least $|S|$
elements of the additive semigroup generated by $S$. The proof combines
Kneser's theorem with explicit representations in residue classes.
The Gaussian binomial result follows from the regular-sequence criterion
of Gatzweiler, Levic\'an-Santib\'a\~nez and Yoshida.
\end{abstract}

\medskip
\noindent\textbf{Keywords.} Gaussian binomial coefficients; coefficient
positivity; additive semigroups; Kneser's theorem; regular sequences.


\section{Introduction and main results}

For an indeterminate $x$, let
\[
 \qbinom{n}{k}_x
 =\prod_{i=1}^{k}\frac{1-x^{n-k+i}}{1-x^i}
 \qquad(0\leq k\leq n)
\]
be the Gaussian binomial coefficient.

Gatzweiler and Krattenthaler \cite[Conjecture~1]{GK} conjectured that
\begin{equation}\label{eq:conjecture}
 \frac{\qbinom{n}{k}_x}{\qbinom{n}{\ell}_x}\in\ZZ[x]
 \quad\Longrightarrow\quad
 \frac{\qbinom{n}{k}_x}{\qbinom{n}{\ell}_x}\in\ZZ_{\geq0}[x].
\end{equation}
After applying the symmetry of Gaussian binomial coefficients, the
nontrivial range is $0\leq\ell<k\leq n/2$. Their paper discusses the
cases $\ell\leq2$ and $k-\ell\leq3$, including proofs for $\ell=2$
and $k-\ell=3$, and reports computational verification for $n\leq400$.

Gatzweiler, Levic\'an-Santib\'a\~nez and Yoshida \cite{GLY} established
a semigroup criterion for a product quotient to be the Hilbert series
of a finite-dimensional graded complete intersection. As an application,
they proved \eqref{eq:conjecture} when
\[
 2k\left\lfloor\frac{k-1}{2}\right\rfloor+k-\ell<n,
\]
and in particular when $n\geq k^2$
\cite[Corollary~3.10]{GLY}. Their criterion asks for a representation
count for every subset of the denominator weights. We prove that the
interval structure of Gaussian binomial quotients supplies all of these
counts, without a restriction on $n$ beyond the original parameter range.

The main result concerns an arbitrary set of distinct generators whose
diameter is smaller than the length of the interval in which representations
are counted. It converts divisibility counts into additive representation
counts. We first state this result, then deduce \eqref{eq:conjecture}.

For a finite set $S$ of positive integers, write
\[
 \Ga(S)=\left\{\sum_{s\in S}n_s s:n_s\in\NN\right\}.
\]
All intervals in this paper are intervals of integers:
$\iv{L}{U}=\{L,L+1,\ldots,U\}$. For a finite set $E$ and a positive
integer $d$, put
\[
 \nu_d(E)=\#\{e\in E:d\mid e\}.
\]
We call an element of $\Ga(S)$ \emph{represented}.

\begin{theorem}[Interval semigroup theorem]\label{thm:interval}
Let $S$ be a nonempty finite set of distinct positive integers, and let
$A=\iv{L}{U}$ have length $r=U-L+1$. Assume that
\begin{equation}\label{eq:admissible}
 L>\max S,\qquad \diam(S)<r,\qquad
 \nu_d(A)\geq\nu_d(S)\quad\text{for every }d\geq1.
\end{equation}
Then
\[
 |A\cap\Ga(S)|\geq |S|.
\]
\end{theorem}

The distinctness hypothesis concerns $S$ itself. In particular, the
theorem makes no assertion with repeated generator weights counted
with multiplicity.

\begin{corollary}\label{cor:GK}
For integers $n\geq0$ and $0\leq k,\ell\leq n$, if
\[
 R_{n,k,\ell}(x)=\frac{\qbinom{n}{k}_x}{\qbinom{n}{\ell}_x}
 \in\ZZ[x],
\]
then $R_{n,k,\ell}(x)\in\ZZ_{\geq0}[x]$.
\end{corollary}

Corollary~\ref{cor:GK} is the assertion formulated by Gatzweiler and
Krattenthaler in \cite[Conjecture~1]{GK}. Its deduction from
Theorem~\ref{thm:interval}, using \cite[Propositions~3.3 and~3.6]{GLY},
is given in Section~\ref{sec:GK}. In this paper $x$ is the polynomial
indeterminate; the integer $q$ appearing in the proof is the quotient
in the division of an interval length by a modulus.

The counting argument below separates two requirements. An existing
generator residue needs only to contribute enough represented points
to meet its count. Residues obtained by adding a deficient residue
must be represented throughout the interval. This distinction allows
the argument to use partial information in the original residue classes.
If the relevant sumset grows through a uniquely represented element,
Kneser's theorem gives enough new classes. In the other case, its
stabilizer is a subgroup, and the original divisor counts control the
points in that subgroup.

Section~2 develops this residue argument. Section~3 records two
elementary facts about two-generator semigroups. The proof of
Theorem~\ref{thm:interval} in Section~4 chooses the smallest generator
as modulus when $r\geq\min S$; otherwise a putative failure supplies
a pair of generators with a common multiple in $A$, whose greatest
common divisor supplies the modulus. Section~\ref{sec:GK} applies the
result to Gaussian binomial quotients. No finite computation is used
in the proof.

\section{Finite sumsets and a residue repair lemma}

For subsets $X,Y$ of an abelian group, write
$X+Y=\{x+y:x\in X,\ y\in Y\}$ and
\[
 \stab(X)=\{h:X+h=X\}.
\]
We use the following form of Kneser's theorem; see
\cite[Theorem~1]{DeVos}.

\begin{theorem}[Kneser]\label{thm:Kneser}
If $X,Y$ are finite nonempty subsets of an abelian group and
$H=\stab(X+Y)$, then
\[
 |X+Y|\geq |X+H|+|Y+H|-|H|.
\]
\end{theorem}

\begin{lemma}[Unique representation]\label{lem:unique}
Let $X,Y$ be finite nonempty subsets of an abelian group. If some
$\gamma\in X+Y$ has exactly one ordered representation
$\gamma=x_0+y_0$ with $x_0\in X$ and $y_0\in Y$, then
\[
 |X+Y|\geq |X|+|Y|-1.
\]
\end{lemma}

\begin{proof}
Put $H=\stab(X+Y)$ and consider the subsets of $H$
\[
 X_0=(X\cap(x_0+H))-x_0,
 \qquad
 Y_0=-(Y\cap(y_0+H))+y_0.
\]
Uniqueness says $X_0\cap Y_0=\{0\}$, so
$|X_0|+|Y_0|\leq |H|+1$. Count missing elements in $x_0+H$ relative
to $X$, and in $y_0+H$ relative to $Y$, separately. Their total is at
least $|H|-1$, even if these two cosets coincide as subsets of the group.
Consequently
\[
 |X+H|-|X|+|Y+H|-|Y|\geq |H|-1.
\]
Substitution in Theorem~\ref{thm:Kneser} proves the claim.
\end{proof}

For the rest of this section suppose \eqref{eq:admissible} holds. Fix
a positive integer $g$ such that $r\geq g$ and at least one element of
$S$ is divisible by $g$. Write
\begin{equation}\label{eq:division}
 r=qg+d,\qquad q\geq1,\quad 0\leq d<g.
\end{equation}
For $\rho\in\ZZ/g\ZZ$, define
\[
 A_\rho=\{x\in A:x\equiv\rho\pmod g\},\qquad
 a_\rho=|A_\rho|,\qquad
 s_\rho=\#\{s\in S:s\equiv\rho\pmod g\}.
\]
Let
\begin{equation}\label{eq:BD}
 B=\{\rho:s_\rho>0\},\qquad
 D=\{\rho:s_\rho>a_\rho\},\qquad
 D_0=\{0\}\cup D.
\end{equation}
Here $B$ denotes a set of residues. We have $0\in B$, and the divisor
condition at $g$ gives $a_0\geq s_0$. Thus
\begin{equation}\label{eq:zero-not-D}
 0\notin D,\qquad |D_0|=|D|+1.
\end{equation}

Every $a_\rho$ is $q$ or $q+1$. The strict diameter condition gives
$s_\rho\leq q+1$; if $d=0$, it gives $s_\rho\leq q$. Hence each
deficient class, namely each member of $D$, has deficit exactly one.
If $D\ne\varnothing$, then $d>0$, and each deficient class has $q+1$
generators and $q$ interval points. Its generators form an exact
progression
\begin{equation}\label{eq:deficientAP}
 w,w+g,\ldots,w+qg.
\end{equation}
Indeed, $q+1$ distinct elements in one residue have span at least $qg$;
omitting any intermediate position would give span at least
$(q+1)g>r$.

We shall say that a residue is \emph{fully represented in $A$} if all
its points in $A$ belong to $\Ga(S)$. If $[y]_g$ denotes the least
nonnegative remainder modulo $g$, then, when $d>0$,
\begin{equation}\label{eq:short-class}
 a_\rho=q\quad\Longleftrightarrow\quad [U-\rho]_g\geq d.
\end{equation}
This follows by viewing the last $d$ positions of the interval as the
extra positions beyond $q$ complete residue systems.

\begin{lemma}[Residue repair]\label{lem:repair}
Retain all assumptions in \eqref{eq:admissible}, the choice of $g$
in \eqref{eq:division}, and the notation \eqref{eq:BD}. Assume that:
\begin{enumerate}
 \item every $\rho\in B$ has at least
 $\min(s_\rho,a_\rho)$ represented points in $A$;
 \item every residue in $B+D$ is fully represented in $A$.
\end{enumerate}
Then $|A\cap\Ga(S)|\geq |S|$.
\end{lemma}

\begin{proof}
The base classes $B$ contribute at least $|S|-|D|$ represented points.
If $D$ is empty this suffices. Suppose $D\ne\varnothing$, and put
\[
 C=B+D,\qquad E=B+D_0=B\cup C.
\]
First suppose $B\subseteq C$. Then $E=C$ is fully
represented. Let $H=\stab(E)$. Since $0\in E$ and $E+H=E$, we have
$H\subseteq E$. Write $H=h\ZZ/g\ZZ$ for its associated divisor
$h$ of $g$. All points of $A$ whose residues lie in $H$ are represented,
and the divisor condition at $h$ gives at least as many such points as
there are generators in $H$.

Outside $H$, the classes $B$ contribute at least their generators'
count minus $|D\setminus H|$. Kneser's theorem gives
\begin{align*}
 |E\setminus(B\cup H)|
 &=|E|-|B\cup H|\\
 &\geq |B+H|+|D_0+H|-|H|-|B\cup H|\\
 &\geq |D_0+H|-|H|\\
 &\geq |D\setminus H|.
\end{align*}
Here $B\cup H\subseteq B+H$ because $0\in B$, and
$D\setminus H\subseteq(D_0+H)\setminus H$ because $0\in D_0$.
Each new residue has at least one point in $A$, since $r\geq g$, and
is fully represented. These disjoint contributions repair all remaining
deficits outside $H$.

Now suppose $\beta\in B\setminus C$. Its only ordered
representation in $E=B+D_0$ is $\beta+0$: any representation with a
summand in $D$ would put $\beta$ in $C$. Lemma~\ref{lem:unique}
therefore gives
\[
 |E|\geq |B|+|D_0|-1=|B|+|D|.
\]
Thus $E\setminus B=C\setminus B$ has at least $|D|$ new
residues. All are fully represented, and each contains at least one
point of $A$. Added to the $|S|-|D|$ base points, they give the desired
count.
\end{proof}

\begin{remark}\label{rem:partial}
The first hypothesis of Lemma~\ref{lem:repair} is a count bound, whereas
the second is a full representation statement. If $B$ is not contained
in $B+D$, unique representation yields enough new residues. If it is
contained, the stabilizer subgroup is fully represented and the original
divisor inequalities repair the deficits in that subgroup. This is the
mechanism that permits partially represented base classes.
\end{remark}

\section{Two elementary semigroup facts}

\begin{lemma}\label{lem:pair}
Let $a<b$ be coprime positive integers, put $\delta=b-a$, and set
$F=ab-a-b$.
\begin{enumerate}
 \item Every integer $n>F$ lies in $\Ga(\{a,b\})$.
 \item For an integer $q\geq\delta-1$, put
 \[
 j_* =\max\left(0,\left\lceil\frac{a-q-1}{\delta}\right\rceil\right).
 \]
 Then $\Ga(\{a,b\})+\iv{0}{q}$ contains every integer at least $aj_*$.
\end{enumerate}
\end{lemma}

\begin{proof}
For the first assertion choose $j\in\{0,\ldots,a-1\}$ such that
$jb\equiv n\pmod a$. Since $n\geq(a-1)(b-1)$,
\[
 n-jb\geq(a-1)(b-1)-(a-1)b=-(a-1).
\]
As a multiple of $a$, this quantity is nonnegative.

For the second assertion, sums of exactly $k$ elements chosen from
$\{a,b\}$ form $ka+\delta\iv{0}{k}$. Adding $\iv{0}{q}$ fills
the whole interval $\iv{ka}{kb+q}$ because $q\geq\delta-1$.
Successive such intervals join without a missing integer whenever
\[
 (k+1)a\leq kb+q+1,
 \quad\text{equivalently}\quad k\delta+q+1\geq a.
\]
This holds for all $k\geq j_*$, so their union contains every integer
from $aj_*$ onward.
\end{proof}

\section{Proof of the interval semigroup theorem}

\begin{proof}[Proof of Theorem~\ref{thm:interval}]
We first treat $r\geq\min S$, and then exclude a failure when
$r<\min S$.

\medskip
\noindent\emph{The case $r\geq\min S$.}
Choose $g=\min S$ and use \eqref{eq:division}--\eqref{eq:BD}. Every
generator residue is fully represented: a point $x\in A$ congruent
to $s\in S$ satisfies $x>s$, and $x-s$ is a nonnegative multiple of
the generator $g$. Thus the first hypothesis of
Lemma~\ref{lem:repair} holds. If $D$ is empty we are done.

Otherwise $d>0$. A deficient class has the progression
\eqref{eq:deficientAP}. Its start satisfies $w>g$, since its residue
is nonzero and $g$ is the smallest generator. Comparing the top element
with $g$ gives $w\leq g+d-1$. Write
\[
 w=g+u,\qquad 1\leq u\leq d-1,\qquad
 s=w+qg=(q+1)g+u.
\]
The divisor condition at the actual generator $s$ supplies a multiple
$ms\in A$. Since $L>s$, we have $m\geq2$ and hence $U\geq2s$.
Set $\beta=U-2s\geq0$. The deficient class is short, so
\eqref{eq:short-class} gives
\[
 [U-u]_g=[u+\beta]_g\geq d,
 \qquad\text{and consequently}\qquad u+\beta\geq d.
\]
This uses only $U\geq2s$; it does not require the double $2s$ itself
to belong to $A$.

Put $M=\max S$. The diameter condition and the expression for $U$ give
\[
 M\leq(q+1)g+d-1,\qquad
 L=(q+2)g+2u+\beta-d+1.
\]
Therefore
\begin{equation}\label{eq:long-threshold}
 L-(M+w-g)\geq g+u+\beta-2d+2\geq g-d+2>0.
\end{equation}
For any generator $z$, every point of its sum residue $z+w\pmod g$
from $z+w$ onward is represented, since $g$ is a generator.
By \eqref{eq:long-threshold}, every point of that residue in $A$ is
greater than $z+w-g$, and hence is at least $z+w$. Thus all classes
of $B+D$ are fully represented. Lemma~\ref{lem:repair} finishes the
long case.

\medskip
\noindent\emph{The short case: choice of modulus.}
Suppose $r<\min S$ and, for contradiction, $|A\cap\Ga(S)|<|S|$.
For each $s\in S$, its divisor condition supplies a positive multiple
in $A$. These chosen multiples belong to fewer than $|S|$ represented
points, so two distinct generators have a common multiple in $A$.
Write them
\[
 ga<gb,\qquad \gcd(a,b)=1,\qquad Tgab\in A\quad(T\geq1).
\]
The strict diameter gives $r>g(b-a)\geq g$. With
$r=qg+d$ as in \eqref{eq:division}, the short interval gives
\begin{equation}\label{eq:short-parameters}
 a\geq q+1,\qquad \delta=b-a\leq q,\qquad b\geq q+2.
\end{equation}
Only $U\geq gab$ will be needed; no assumption on the size of $T$
is imposed.

Set $F=ab-a-b$. Every $g$-multiple in $A$ is represented by the pair,
since
\begin{equation}\label{eq:zero-full}
 L\geq Tgab-r+1>g(ab-a)>gF.
\end{equation}
The divisor condition at $g$ also supplies at least the number of
generators in the zero class.

\medskip
\noindent\emph{Base counts in nonzero residues.}
Consider a nonzero residue containing $c$ generators and let $w$ be its
smallest generator. Its largest is at least $w+(c-1)g$. Comparing that
element with $ga$ gives
\begin{equation}\label{eq:base-start}
 w\leq ga+(q-c+1)g+d-1.
\end{equation}
The last point of this residue not exceeding $U$ is $w+Ng$, where
\begin{align*}
 N&=\left\lfloor\frac{U-w}{g}\right\rfloor\\
 &\geq ab-a-q+c-1+\left\lfloor\frac{1-d}{g}\right\rfloor
 \geq ab-a-q+c-2.
\end{align*}
The first inequality uses $U\geq gab$ and \eqref{eq:base-start}.
The second follows from $0\leq d\leq g-1$, which implies
$\lfloor(1-d)/g\rfloor\geq-1$. By \eqref{eq:short-parameters},
\begin{equation}\label{eq:base-tail}
 N-F\geq b-q+c-2\geq c.
\end{equation}
The last $\min(c,a_\rho)$ points of this residue in $A$ therefore have
offsets at least $N-c+1\geq F+1$. They are generated by $w$ together
with $ga,gb$. Along with \eqref{eq:zero-full}, this verifies the first
hypothesis of Lemma~\ref{lem:repair}.

If $D$ is empty, that hypothesis already gives $|S|$ points. Otherwise
$d>0$, and each deficient class is the progression
\begin{equation}\label{eq:shortAP}
 w,w+g,\ldots,w+qg,\qquad w\leq ga+d-1.
\end{equation}
Its $q$ points in $A$ are all represented by \eqref{eq:base-tail}.
It remains to fill the sum residues $B+D$.

\medskip
\noindent\emph{Sum residues when $q\geq2$.}
Let $s$ be any actual generator and $w$ a deficient AP start as in
\eqref{eq:shortAP}. Adding $s$ to that progression and applying
Lemma~\ref{lem:pair}(2) represents every point of the sum residue
from
\[
 V=s+w+ga j_*,\qquad
 j_* =\max\left(0,\left\lceil\frac{a-q-1}{\delta}\right\rceil\right),
\]
onward, in steps of $g$. These are actual semigroup sums, and using
the same generator more than once is allowed.

The diameter condition gives
$s\leq\max S\leq ga+qg+d-1$. Together with
\eqref{eq:shortAP}, this yields
\[
 V-g\leq g(2a+q+aj_*-1)+2d-2.
\]
Also $L\geq gab-qg-d+1$. Since $a\geq q+1$,
$j_*\leq a-q-1$. Consequently
\begin{align}
 L-(V-g)
 &\geq g[ab-2a-aj_*-2q+1]-3d+3\notag\\
 &\geq g[a(\delta-1)+q(a-2)+1]-3d+3\notag\\
 &\geq 3g-3d+3\geq6.\label{eq:q-large-threshold}
\end{align}
The final line uses $q\geq2$, $a\geq q+1$, and $d\leq g-1$.
Thus every point of the sum residue in $A$ is greater than $V-g$,
and hence is at least $V$. All classes in $B+D$ are fully represented,
and Lemma~\ref{lem:repair} gives the contradiction.

\medskip
\noindent\emph{Sum residues when $q=1$.}
Now $\delta=1$, $b=a+1$, $a\geq2$, and $j_*=a-2$. Write a deficient
AP start as
\[
 w=ga+u,\qquad u\leq d-1<g.
\]
The integer $u$ may be negative; no lower bound $u\geq0$ is assumed.
The AP is $w,w+g$. For every generator $s$, the same threshold satisfies
\begin{equation}\label{eq:q-one-upper}
 V-g=s+w+ga(a-2)-g\leq ga^2+u+d-1,
\end{equation}
because $\max S\leq(a+1)g+d-1$.

If $a\geq3$, write
\[
 U=ga(a+1)+\beta,\qquad\beta\geq0.
\]
The deficient class has exactly one point of $A$, so its short-class
condition is
\[
 [U-w]_g=[\beta-u]_g\geq d.
\]
Since $\beta-u>-g$, this implies $\beta-u\geq d-g$. Indeed,
if $\beta-u<0$, its least nonnegative remainder is $\beta-u+g$;
if $\beta-u\geq0$, the bound follows from $d-g<0$. Thus
\begin{align}
 L-(ga^2+u+d-1)
 &=g(a-1)+\beta-u-2d+2\notag\\
 &\geq g(a-2)-d+2>0.\label{eq:q-one-a-large}
\end{align}

If $a=2$, write $U=6g+\beta$ with $\beta\geq0$. The top AP generator
is $w+g=3g+u$. Its supplied multiple in $A$ is at least twice this
generator, whence $\beta\geq2u$. Together with $\beta\geq0$, this
gives $\beta-u\geq0$, including when $u$ is negative. The short-class
condition $[\beta-u]_g\geq d$ now implies $\beta-u\geq d$. Since
$L=5g+\beta-d+1$, we obtain
\begin{equation}\label{eq:q-one-a-two}
 L-(4g+u+d-1)=g+\beta-u-2d+2\geq g-d+2>0.
\end{equation}
Equations \eqref{eq:q-one-upper}--\eqref{eq:q-one-a-two} show that
$L>V-g$ in both cases. Hence every required sum residue is fully
represented in $A$, and Lemma~\ref{lem:repair} again contradicts
failure.

The long and short cases exhaust all possibilities, proving the theorem.
\end{proof}

\section{Gaussian binomial positivity}\label{sec:GK}

We use the following consequence of the regular-sequence criterion of
Gatzweiler, Levic\'an-Santib\'a\~nez and Yoshida
\cite[Propositions~3.3 and~3.6]{GLY}.

\begin{proposition}[Semigroup criterion for positivity]\label{prop:GLY}
Let $a_1,\ldots,a_m,b_1,\ldots,b_m$ be positive integers. If, for every
subset $I\subseteq\{1,\ldots,m\}$,
\[
 \#\{j:a_j\in\Ga(\{b_i:i\in I\})\}\geq |I|,
\]
then
\[
 \frac{\prod_{j=1}^m(1-x^{a_j})}{\prod_{i=1}^m(1-x^{b_i})}
 \in\ZZ_{\geq0}[x].
\]
\end{proposition}

Their criterion supplies a homogeneous regular sequence of degrees
$a_1,\ldots,a_m$ in a polynomial ring with variable weights
$b_1,\ldots,b_m$. The displayed quotient is the Hilbert series of its
finite-dimensional graded quotient, so its coefficients are nonnegative
integer dimensions. This application uses the square case with equally
many degrees and weights.

\begin{proof}[Proof of Corollary~\ref{cor:GK}]
Use the symmetry $\qbinom{n}{j}_x=\qbinom{n}{n-j}_x$ independently in the
numerator and denominator. We may therefore assume
$0\leq k,\ell\leq\lfloor n/2\rfloor$. The Gaussian polynomial has
degree $j(n-j)$, which is strictly increasing on these integer arguments.
For $0\leq u<v\leq\lfloor n/2\rfloor$, the degree difference is
$(v-u)(n-u-v)>0$. Since the quotient is a nonzero polynomial, its
degree is nonnegative; thus $k\geq\ell$. If $k=\ell$, the quotient
is $1$. This also covers $n=0$ and $n=1$.

Suppose $0\leq\ell<k\leq n/2$, and put
\[
 r=k-\ell,\qquad
 A=\iv{n-k+1}{n-\ell},\qquad
 W=\iv{\ell+1}{k}.
\]
Cancellation of Gaussian factorials gives
\begin{equation}\label{eq:GK-product}
 R_{n,k,\ell}(x)
 =\frac{\prod_{a\in A}(1-x^a)}{\prod_{w\in W}(1-x^w)}.
\end{equation}
For every positive integer $d$, the cyclotomic exponent in this quotient
is
\begin{align*}
 \nu_d(A)-\nu_d(W)
 &=\left\lfloor\frac{n-\ell}{d}\right\rfloor
   -\left\lfloor\frac{n-k}{d}\right\rfloor
   -\left\lfloor\frac{k}{d}\right\rfloor
   +\left\lfloor\frac{\ell}{d}\right\rfloor.
\end{align*}
Polynomiality and the distinct irreducibility of cyclotomic polynomials
make these exponents nonnegative for $d\geq2$. For $d=1$, the difference
is zero because $|A|=|W|=r$; for $d>n$, both counts are zero. Hence
\begin{equation}\label{eq:GK-counts}
 \nu_d(A)\geq\nu_d(W)\qquad(d\geq1).
\end{equation}

Let $S$ be any nonempty subset of $W$. We have
\[
 \min A=n-k+1\geq k+1>\max S,
 \qquad \diam(S)\leq k-\ell-1=r-1.
\]
Moreover, \eqref{eq:GK-counts} implies
$\nu_d(A)\geq\nu_d(S)$ for every $d$. Theorem~\ref{thm:interval}
therefore gives
\[
 |A\cap\Ga(S)|\geq |S|
 \qquad\text{for every nonempty }S\subseteq W.
\]
The empty subset imposes only the trivial inequality with right side
zero. Since the weights in $W$ are distinct, these are exactly the
subset conditions in Proposition~\ref{prop:GLY}. Applying that proposition
to \eqref{eq:GK-product} proves
$R_{n,k,\ell}(x)\in\ZZ_{\geq0}[x]$.
\end{proof}

\section*{Acknowledgements and AI assistance}
The author contributed mathematical ideas and proof strategies and used
OpenAI's Codex (GPT-6) in developing and refining the arguments, carrying
out exploratory computations, preparing the manuscript, and checking
consistency. The author takes responsibility for the content of this paper.

\begin{thebibliography}{9}

\bibitem{DeVos}
M.~DeVos,
\emph{A short proof of Kneser's addition theorem for abelian groups},
in M.~B.~Nathanson (ed.), \emph{Combinatorial and Additive Number Theory},
Springer Proceedings in Mathematics \& Statistics \textbf{101},
Springer, 2014, pp.~39--41.
\url{https://doi.org/10.1007/978-1-4939-1601-6_3}.
Also available as arXiv:1303.3539,
\url{https://arxiv.org/abs/1303.3539}.

\bibitem{GLY}
M.~Gatzweiler, F.~Levic\'an-Santib\'a\~nez, and A.~Yoshida,
\emph{Cyclotomic generating functions, empty weighted complete
intersections and positivity},
arXiv:2603.22226v1 (2026).
\url{https://arxiv.org/abs/2603.22226v1}.

\bibitem{GK}
M.~Gatzweiler and C.~Krattenthaler,
\emph{A positivity conjecture for a quotient of $q$-binomial coefficients},
\emph{The Ramanujan Journal} \textbf{69} (2026), Article 13.
\url{https://doi.org/10.1007/s11139-025-01285-2}.
Also available as arXiv:2502.06032,
\url{https://arxiv.org/abs/2502.06032}.

\end{thebibliography}
\end{document}
